\chapter{المقاسات على حلقة مثاليات رئيسية}\label{ch:b3:modules}

جبر خطي على \emph{حلقة} بدل حقل: هذا التغيير الطفيف في الفرضيات
يُنتج واحدة من مبرهنات التوحيد الكبرى في الجبر. \omterm{def:b3:modules:module}{فالمقاس} على $\Z$
زمرة أبيلية؛ \omterm{def:b3:modules:module}{والمقاس} على $K[X]$ فضاء متجهي مزوَّد بتشاكل ذاتي. ومن
ثَمّ فإن مبرهنة البنية للمقاسات المولَّدة بعدد منتهٍ من العناصر على
\omterm{def:b3:rings:pidufd}{حلقة مثاليات رئيسية} تصنّف، دفعةً واحدة، جميع الزمر الأبيلية
المولَّدة بعدد منتهٍ \emph{و} جميع التشاكلات الذاتية إلى حدود
التشابه --- فاختزال جوردان، الذي حصلت عليه السنة الجامعية 2
بتراجعات دقيقة، يسقط نتيجةً فرعية، ومعه شقيقه الأدق: الصيغة
الناظمية \emph{الناطقة}، الصحيحة على كل حقل. والمحرّك الحسابي هو
صيغة سميث الناظمية، وهي حسابٌ للمصفوفات جدير بإقليدس.

في كل ما يلي، $A$ حلقة تبديلية، وستصير قريبًا \omterm{def:b3:rings:pidufd}{حلقة مثاليات رئيسية}؛
وتعني كلمة «\omterm{def:b3:modules:module}{مقاس}» مقاسًا على $A$.

\section{المقاسات والمقاسات الحرة}

\begin{definition}\label{def:b3:modules:module}
\emph{المقاس على $A$}\index{مقاس} هو زمرة أبيلية $(M, +)$ مزوَّدة
بضرب سلَّمي $A \times M \to M$ يحقق بديهيات الفضاءات المتجهية: $a(x + y) = ax + ay$، $(a + b)x = ax + bx$،
$(ab)x = a(bx)$، $1x = x$. أما المقاسات الجزئية، وحواصل القسمة
$M/N$، والتشاكلات (التطبيقات الخطية على $A$)، والمجاميع المباشرة
$\bigoplus_i M_i$، \omterm{thm:b3:groups:firstiso}{ومبرهنات التماثل}، فتُعرَّف ويُبرهَن عليها حرفيًا كما في حالة
الفضاءات المتجهية والزمر الأبيلية؛ وعلى وجه الخصوص $M/\ker f \cong
\operatorname{im} f$ من أجل
تشاكل $f$.
\end{definition}

\begin{example}\label{ex:b3:modules:examples}
الحالات الثلاث المحرّكة.
\begin{enumerate}
  \item $A = K$ حقل: فالمقاسات فضاءات متجهية.
  \item $A = \Z$: فالمقاسات هي بالضبط الزمر الأبيلية (إذ يُفرض
        على $nx$ أن يكون $x + \dots + x$)، والمقاسات الجزئية هي
        الزمر الجزئية.
  \item $A = K[X]$: \omterm{def:b3:modules:module}{فالمقاس} فضاء متجهي $V$ على $K$ مزوَّد بتطبيق
        خطي على $K$ هو $u\colon x \mapsto X\cdot x$ --- وبالعكس، يصير كل زوج $(V, u)$
        مع $u \in \mathcal L(V)$ مقاسًا على $K[X]$ بواسطة $P \cdot x = P(u)(x)$. والمقاسات
        الجزئية هي بالضبط الفضاءات الجزئية الصامدة تحت $u$.
\end{enumerate}
\omterm{def:b3:rings:ideal}{والمثالي} في $A$ هو بالضبط \omterm{def:b3:modules:module}{مقاس} جزئي من $A$؛ \omterm{def:b3:rings:ideal}{وحلقة القسمة} $A/I$
مقاسٌ على $A$. وخلافًا للفضاءات المتجهية، يمكن أن يكون للمقاسات
\emph{\omterm{def:b3:modules:torsion}{التواء}}: ففي $\Z/6\Z$، يُبيد العنصرَ $\bar 3 \ne 0$ العنصرُ
$2 \neq 0$.
\end{example}

\begin{definition}\label{def:b3:modules:free}
يكون $M$ \emph{مولَّدًا بعدد منتهٍ من العناصر}\index{مقاس مولد بعدد
منته} إذا كان $M = Ax_1 + \dots + Ax_n$ من أجل عناصر $x_i$ ما. ويكون $M$
\emph{حرًّا}\index{مقاس حر} من الرتبة $n$ إذا كان $M \cong A^n$، أي
إذا كان له \emph{أساس} (عائلة مولِّدة ومستقلة خطيًا على $A$). وكل
\omterm{def:b3:modules:module}{مقاس} $M$ مولَّد بعدد منتهٍ هو حاصل قسمة \omterm{def:b3:modules:module}{لمقاس} حر: إذ يرسل $(a_1, \dots, a_n) \mapsto \sum a_ix_i$
المقاسَ $A^n$ على $M$.
\end{definition}

\begin{proposition}[صمود الرتبة]\label{prop:b3:modules:rank}
إذا كان $A \neq 0$ و $A^m \cong A^n$، فإن $m = n$.
\end{proposition}

\begin{proof}
نختار مثاليًا أعظميًا $\mathfrak m$ في $A$ (\cref{thm:b3:rings:krull}) ونضع
$k = A/\mathfrak m$، وهو حقل. ويرسل التماثل $f \colon A^m
\to A^n$ \omterm{def:b3:modules:module}{المقاس} $\mathfrak m A^m$ داخل $\mathfrak m A^n$
(بالخطية)، ومنه فهو يُحدث تماثلًا بين حاصلَي القسمة
\[
A^m/\mathfrak m A^m \;\cong\; A^n/\mathfrak m A^n,
\qquad\text{أي}\qquad k^m \cong k^n
\]
بوصفهما فضاءين متجهيين على $k$ (إذ يُبيد $\mathfrak m$ حاصلَ
القسمة $A^m/\mathfrak m A^m$، ومنه فإن فعل $A$ يتحلّل عبر $k$؛ وصور الأساس
القانوني تشكّل أساسًا على $k$). وتعطي نظرية البُعد على الحقل $k$
أن $m = n$.
\end{proof}

\begin{theorem}[المقاسات الجزئية من مقاس حر]\label{thm:b3:modules:submodule}
لتكن $A$ \omterm{def:b3:rings:pidufd}{حلقة مثاليات رئيسية} وليكن $M \subseteq A^n$ مقاسًا جزئيًا. عندئذٍ
يكون $M$ \omterm{def:b3:modules:free}{حرًّا} ورتبته $\leq n$.
\end{theorem}

\begin{proof}
بالتراجع على $n$. من أجل $n = 1$: يكون $M$ مثاليًا، ومنه $M = (0)$
(وهو حر من الرتبة $0$) أو $M = dA \cong A$ (إذ إن $x \mapsto dx$
متباين لأن الحلقة تامة). ومن أجل $n > 1$: ليكن $\pi \colon A^n \to A$ الإحداثي
الأخير. عندئذٍ يكون $\pi(M)$ مثاليًا، أي $(0)$ أو $dA$. ففي حالة
$(0)$: يكون $M \subseteq A^{n-1} \times \{0\}$ وينطبق التراجع. وإلا فلنختر $x_0 \in M$ يحقق
$\pi(x_0) = d$. عندئذٍ يُكتب كل $x \in M$ بطريقة وحيدة
\[
x = \underbrace{\frac{\pi(x)}{d}}_{\in A}\, x_0 +
\Bigl(x - \tfrac{\pi(x)}d x_0\Bigr),
\qquad x - \tfrac{\pi(x)}d x_0 \in M \cap \ker\pi
\]
(لأن $\pi(x) \in dA$، ومنه فالمعامل في $A$). ومنه $M = Ax_0
\oplus (M \cap \ker \pi)$: والمجموع
مباشر لأن $\pi(ax_0) =
ad = 0$ يفرض $a = 0$. وبالتراجع يكون $M \cap \ker\pi \subseteq
\ker \pi \cong A^{n-1}$ \omterm{def:b3:modules:free}{حرًّا} ورتبته
$\leq n - 1$؛ وبإضافة $x_0$ (المستقل عن $\ker\pi$ كما رأينا للتوّ)
نحصل على أساس \omterm{def:b3:modules:module}{للمقاس} $M$ عدد عناصره $\leq n$.
\end{proof}

\begin{remark}\label{rem:b3:modules:presentation}
ومن نتائج ذلك أن لكل \omterm{def:b3:modules:module}{مقاس} $M$ مولَّد بعدد منتهٍ على \omterm{def:b3:rings:pidufd}{حلقة مثاليات
رئيسية} \emph{تقديمًا منتهيًا}: فلتطبيق شامل $\varphi\colon A^n \to M$ نواةٌ حرة أساسها
$c_1, \dots, c_k$ (حيث $k \leq n$)، ويكون $M
\cong A^n / \operatorname{im}(C)$ حيث $C \in M_{n,k}(A)$ هي المصفوفة التي
أعمدتها العناصر $c_j$. ففهم $M$ يعني فهم \emph{مصفوفة على $A$}
إلى حدود تغيير الأسس في المنطلق والمستقر --- وهو موضوع القسم
التالي.
\end{remark}

\section{صيغة سميث الناظمية}

\begin{definition}\label{def:b3:modules:equivalence}
نقول عن مصفوفتين $B, C \in M_{n,k}(A)$ إنهما \emph{متكافئتان} إذا كان $C =
QBP$ حيث
$Q \in GL_n(A)$ و $P \in GL_k(A)$ (أي قابلتان للقلب \emph{على
$A$}: بمحدد في $A^\times$). ومصفوفتا تقديم متكافئتان تعرّفان
مقاسين متماثلين $A^n/
\operatorname{im}$ (بتغيير الأسس في $A^n$ و $A^k$).
\end{definition}

\begin{theorem}[صيغة سميث الناظمية]\label{thm:b3:modules:smith}
لتكن $A$ \omterm{def:b3:rings:pidufd}{حلقة مثاليات رئيسية} ولتكن $B \in M_{n,k}(A)$. عندئذٍ تكافئ $B$ مصفوفة
قطرية\index{صيغة سميث الناظمية}
\[
\operatorname{diag}(d_1, d_2, \dots, d_r, 0, \dots, 0),
\qquad d_1 \mid d_2 \mid \cdots \mid d_r \neq 0 ,
\]
وتكون العناصر $d_i$ وحيدة إلى حدود التصاحب: إذ إن
$d_1 \cdots d_i$ \omterm{lem:b3:rings:bezout}{قاسم مشترك أكبر} لمحدّدات $B$ الجزئية من القياس
$i \times i$ (وعلى وجه الخصوص فإن هذا القاسم صامد بالتكافؤ). وتُسمّى
العناصر $d_i$ \emph{العوامل الصامدة}\index{عوامل صامدة} للمصفوفة
$B$.
\end{theorem}

\begin{proof}
\emph{الوجود.} إذا كان $B = 0$ فقد انتهينا. وإلا فلننظر في مجموعة
المثاليات $(b)$ المولَّدة بعناصر مصفوفات تكافئ $B$؛ وبما أن $A$
\omterm{def:b3:rings:noetherian}{نويثرية}، فلنختر مصفوفة $B'$ تكافئ $B$ وعنصرًا $d$ فيها بحيث يكون
$(d)$ \emph{أعظميًا} في هذه المجموعة. وننقل $d$ إلى الموضع $(1,1)$
بتبديل الأسطر والأعمدة.

\emph{الادعاء: يقسم $d$ كل عناصر $B'$.} أولًا العمود 1: إذا لم يكن
$b_{i1}$ مضاعفًا للعنصر $d$، فليكن $e = \gcd(d, b_{i1}) = ud +
vb_{i1}$ (بيزو)، ومنه $(e) \supsetneq (d)$. وحيلة
المصفوفة من القياس $2 \times 2$: بالفعل على السطرين $1$ و $i$
بواسطة
\[
\begin{pmatrix} u & v \\ -\dfrac{b_{i1}}e & \dfrac de
\end{pmatrix}
\in GL_2(A)
\qquad
\Bigl(\det = \tfrac{ud + v b_{i1}}e = 1\Bigr)
\]
نحصل على مصفوفة مكافئة عنصرها في الموضع $(1,1)$ هو $e$: وهذا
يناقض أعظمية $(d)$. إذن يقسم $d$ العمودَ $1$، وبالتناظر السطرَ
$1$. وبطرح مضاعفات السطر 1 والعمود 1 نُصفّرهما: فتكافئ $B'$ المصفوفة
$\begin{pmatrix} d &
0\\ 0 & B''\end{pmatrix}$. وبعد ذلك، يقسم $d$ كل عنصر $b$ من $B''$: نضيف سطر $b$
إلى السطر 1 (وهي عملية أوّلية؛ فيحتوي السطر الأول الجديد على $d$
وعلى عناصر $b$)، ونكرّر حجّة تصفير العمود: إذ إن عنصرًا غير مضاعف
سيحسّن $(d)$ مرة أخرى. والآن نتراجع على القياس: فللمصفوفة $B''$، التي
تقبل جميع عناصرها القسمة على $d$، صيغةُ سميث $\operatorname{diag}(d_2,
\dots)$ تبقى عناصرها
قابلة للقسمة على $d$ (لأن كل عنصر من أي $QB''P$ تركيبةٌ على $A$
لعناصر $B''$)؛ ونضع $d_1 =
d$.

\emph{الوحدانية.} ليرمز $D_i(B)$ إلى \omterm{lem:b3:rings:bezout}{قاسم مشترك أكبر} لجميع
المحدّدات الجزئية من القياس $i \times i$. لا يمكن للعمليات على
الأسطر والأعمدة، ولا بصورة أعمّ للضرب في \emph{أي} مصفوفة، أن
تصغّر هذا القاسم: فالمحدّدات الجزئية من القياس $i \times i$
للمصفوفة $QB$ تركيباتٌ على $A$ لمثيلاتها في $B$ (بنشر كوشي--بينيه؛
أو مباشرةً: كل سطر من $QB$ تركيبة لأسطر $B$، والمحدّدات الجزئية
متعددة الخطية في الأسطر). ومنه فإن $D_i(QBP)$ و $D_i(B)$ يقسم كل
منهما الآخر: أي إن $D_i$ صامد بالتكافؤ. وعلى الصيغة القطرية، تكون
المحدّدات الجزئية غير المعدومة من القياس $i
\times i$ هي جداءات $i$ من
العناصر $d_j$، وتجعل القابليةُ للقسمة $d_1 \mid \dots \mid d_r$ المقدارَ
$d_1 \cdots d_i$ هو \omterm{lem:b3:rings:bezout}{القاسم المشترك الأكبر}. ومنه $d_1 \cdots d_i = D_i(B)$ إلى حدود
العناصر القابلة للقلب، ويتحدّد $d_i =
D_i/D_{i-1}$.
\end{proof}

\begin{method}\label{met:b3:modules:smith}
على \omterm{def:b3:rings:pidufd}{حلقة إقليدية} ($\Z$، $K[X]$)، يكون اختزال سميث خوارزمية ---
دون حاجة إلى حجّة الأعظمية: ننقل العنصر ذا القياس الإقليدي الأصغر
إلى الموضع $(1,1)$؛ فإذا أخفق في قسمة عنصر ما من سطره أو عموده،
تركت قسمةٌ إقليدية باقيًا أصغر تمامًا هناك --- نبدّله إلى الموضع
ونعيد الكرّة (والانتهاء مضمون: لأن الأقياس تتناقص)؛ وحين يقسم
سطره وعموده بأكملهما، نُصفّرهما؛ وإذا أخفق في قسمة عنصر داخلي،
أضفنا ذلك السطر إلى السطر $1$ وأعدنا الكرّة؛ ثم نتراجع على الكتلة
الداخلية. وعمليًا على المصفوفات الصحيحة: نحسب $D_1 = \gcd$
لعناصرها، ثم $D_2$، \dots\ بالمحدّدات الجزئية من أجل الأقياس
الصغيرة، أو نشغّل الخوارزمية.
\end{method}

\begin{example}[اختزال سميث كاملًا]
\label{ex:b3:modules:smithworked}
لنختزل $M = \begin{pmatrix} 1 & 2 & 3\\ 4 & 5 & 6\\ 7 & 8 & 9
\end{pmatrix}$ على $\Z$. الركن $1$ يقسم كل شيء: فنُصفّر سطره
وعموده ($L_2 \leftarrow L_2 - 4L_1$، $L_3
\leftarrow L_3 - 7L_1$، ثم $C_2 \leftarrow C_2 - 2C_1$، $C_3
\leftarrow C_3 - 3C_1$):
\[
M \sim \begin{pmatrix} 1 & 0 & 0\\ 0 & -3 & -6\\ 0 & -6 & -12
\end{pmatrix} .
\]
وفي الكتلة الداخلية، يقسم الركن $-3$ جميع العناصر: فيصفّرها
$L_3
\leftarrow L_3 - 2L_2$ و $C_3 \leftarrow C_3 - 2C_2$ لتصير $\operatorname{diag}(-3, 0)$. وبضبط الإشارات (بضرب سطر في $-1$،
وهي عملية مشروعة):
\[
M \sim \operatorname{diag}(1, 3, 0),
\qquad
\Z^3/M\Z^3 \cong \Z/3\Z \times \Z .
\]
ولنتحقق بواسطة القواسم المحدّدية: $D_1 =
\gcd(\text{العناصر}) = 1$؛ وكل محدد جزئي من
القياس $2\times2$ للمصفوفة $M$ مضاعفٌ للعدد $3$ (مثلًا $\det\bigl(\begin{smallmatrix}1 & 2\\ 4
& 5\end{smallmatrix}\bigr) = -3$)
وأحدها يساوي $-3$: ومنه $D_2 =
3$؛ و $D_3 = \det M = 0$. إذن $d_1 = 1$
و $d_2 = 3$ و $d_3 = 0$: أي الجواب نفسه. وهنا درسان: العامل
الصامد المعدوم يسجّل هبوط الرتبة (إذ يكتسب حاصل القسمة حدًّا \omterm{def:b3:modules:free}{حرًّا}
من $\Z$)، وسلسلة القابلية للقسمة $1 \mid 3 \mid 0$ هي شهادة سميث --- فاختزالٌ
قطري يخالف السلسلة (مثل $\operatorname{diag}(2, 3)$، وهو ما قد يُنتجه غير المتأني من
$\bigl(\begin{smallmatrix}2 & 0\\ 0 &
3\end{smallmatrix}\bigr)$ بالتوقف مبكرًا: والصيغة الصحيحة لسميث هي $\operatorname{diag}(1, 6)$، لأن
$D_1 = 1$ هنا!) ليس اختزالًا منتهيًا.
\end{example}

\begin{omfigure}
\begin{tikzpicture}[scale=0.72]
  \clip (-3.4,-1.3) rectangle (4.4,4.3);
  \foreach \x in {-3,...,4}
    \foreach \y in {-1,...,4}
      \fill[black!40] (\x,\y) circle (1.6pt);
  \fill[omDef!12, opacity=0.7] (0,0) -- (1,3) -- (3,3) -- (2,0)
    -- cycle;
  \foreach \a/\b in {0/0, 2/0, 4/0, -2/0, 1/3, 3/3, -1/3, -3/3,
                     2/6, -2/-6, 0/6}
    \fill[omThm] (\a,\b) circle (2.6pt);
  \draw[omThm, thick, ->] (0,0) -- (2,0)
    node[midway, below, font=\small] {$(2,0)$};
  \draw[omThm, thick, ->] (0,0) -- (1,3)
    node[midway, left, font=\small] {$(1,3)$};
\end{tikzpicture}

{\small الشبكة الجزئية $L = \Z(2,0) + \Z(1,3)$ من $\Z^2$ (النقط الحمراء). وصيغة
سميث الناظمية للمصفوفة $\bigl(\begin{smallmatrix} 2 &
1\\ 0 & 3\end{smallmatrix}\bigr)$ هي $\operatorname{diag}(1, 6)$: ففي الأساس الملائم
$f_1 = (1,3)$ و $f_2 = (0,1)$ من $\Z^2$، لدينا $L = \Z f_1 \oplus \Z\,6f_2$، ومنه
$\Z^2/L \cong \Z/6\Z$ --- فالدليل يساوي $\abs{\det} = 6$، وهو مساحة المجال
الأساسي المظلَّل.}
\end{omfigure}

\section{مبرهنة البنية}

\begin{definition}\label{def:b3:modules:torsion}
لتكن $A$ حلقة تامة وليكن $M$ مقاسًا على $A$. \omterm{def:b3:modules:module}{المقاس} الجزئي
\emph{الالتوائي}\index{التواء} هو
\[
T(M) = \{x \in M : ax = 0 \text{ من أجل } a \neq 0\}
\]
(وهو \omterm{def:b3:modules:module}{مقاس} جزئي: إذ إن $ax = by = 0$ يعطي $ab(x + y) = 0$ مع
$ab \ne 0$). ويكون $M$ \emph{عديم الالتواء} إذا كان $T(M) = 0$،
و\emph{مقاسًا التوائيًا} إذا كان $T(M) = M$.
\end{definition}

\begin{theorem}[بنية المقاسات المولَّدة بعدد منتهٍ على حلقة مثاليات
رئيسية]\label{thm:b3:modules:structure}
لتكن $A$ \omterm{def:b3:rings:pidufd}{حلقة مثاليات رئيسية} وليكن $M$ مقاسًا على $A$ \omterm{def:b3:modules:free}{مولَّدًا بعدد
منتهٍ من العناصر}. يوجد عدد وحيد $r \in \N$ وعناصر غير معدومة وغير
قابلة للقلب $d_1 \mid d_2 \mid
\cdots \mid d_s$، وحيدة إلى حدود التصاحب، بحيث\index{مبرهنة
البنية للمقاسات}
\[
M \;\cong\; A^r \,\oplus\, A/(d_1) \oplus \cdots \oplus A/(d_s).
\]
وعلاوة على ذلك $T(M) \cong A/(d_1)\oplus\dots\oplus A/(d_s)$ و $M/T(M)
\cong A^r$: فكل \omterm{def:b3:modules:module}{مقاس} \emph{عديم \omterm{def:b3:modules:torsion}{الالتواء}}
مولَّد بعدد منتهٍ على \omterm{def:b3:rings:pidufd}{حلقة مثاليات رئيسية} هو \omterm{def:b3:modules:free}{مقاس حر}.
\end{theorem}

\begin{proof}
\emph{الوجود.} نقدّم $M \cong A^n/\operatorname{im}(C)$ (\cref{rem:b3:modules:presentation}) ونضع $C$ في صيغة سميث:
فبعد تغييرَي الأساس نجد $M \cong A^n / (d_1A \times \dots
\times d_rA \times 0 \times \dots \times 0) \cong A/(d_1) \oplus
\dots \oplus A/(d_r) \oplus A^{\,n - r}$. ونحذف العوامل التي يكون فيها
$d_i$ قابلًا للقلب (لأن $A/(d_i) = 0$)؛ وتبقى سلسلة القابلية
للقسمة قائمة.

\emph{تعيين \omterm{def:b3:modules:torsion}{الالتواء}.} في هذا التفكيك، يكون $A^r$ عديم \omterm{def:b3:modules:torsion}{الالتواء}
(لأنه ليس في حلقة تامة قواسم للصفر) ويكون كل $A/(d_i)$ التوائيًا
(لأن $d_i \ne 0$ يبيده)؛ ويشطر المجموع المباشر الالتواءَ تبعًا
لذلك: أي $T(M) = \bigoplus_i A/(d_i)$ و $M/T(M) \cong A^r$.

\emph{وحدانية $r$}: إن $M/T(M) \cong A^r$ يتعلق \omterm{def:b3:modules:module}{بالمقاس} $M$ وحده، وتحدّد
\cref{prop:b3:modules:rank} العددَ $r$.

\emph{وحدانية العناصر $d_i$}: يكفي أن نعالج \omterm{def:b3:modules:module}{المقاس} \omterm{def:b3:modules:torsion}{الالتوائي}
$T = T(M)$. نفكّك كل $d_i$ إلى عناصر أولية ونشطر بمبرهنة الباقي
الصيني (\cref{thm:b3:rings:crt}؛ إذ إن العناصر الأولية المتمايزة تولّد \omterm{thm:b3:rings:crt}{مثاليات
متكاملة}):
\[
A/(d) \cong \bigoplus_{p \mid d} A/\bigl(p^{v_p(d)}\bigr):
\qquad
T \cong \bigoplus_{p} \bigoplus_{j} A/\bigl(p^{k_{p,j}}\bigr),
\]
وتُسمّى العناصر $p^{k_{p,j}}$ \emph{القواسم
الأوّلية}\index{قواسم أولية}. وبالعكس، تُستعاد العناصر $d_i$ من
المجموعة المتعدّدة للقواسم الأوّلية (إذ إن $d_s$ = جداء أعلى قوة
لكل \omterm{def:b3:rings:divisibility}{عنصر أولي}، وهكذا)، ومنه يكفي أن نبرهن على أن المجموعة
المتعدّدة $\{k_{p,j}\}_j$ يحدّدها $T$، من أجل كل \omterm{def:b3:rings:divisibility}{عنصر أولي} $p$.
لنثبّت $p$؛ ومن أجل $j \geq 1$ لننظر في الفضاءات المتجهية على
$A/(p)$ المعطاة بالحواصل $p^{j-1}T/p^jT$. فعلى عامل دائري $A/(p^k)$:
\[
p^{j-1}\bigl(A/(p^k)\bigr)\big/p^{j}\bigl(A/(p^k)\bigr)
\cong
\begin{cases}
A/(p) & \text{إذا كان } j \leq k,\\
0 & \text{إذا كان } j > k,
\end{cases}
\]
وعلى عامل $A/(q^k)$ حيث $q \neq p$: يكون الضرب في $p$ تقابليًا
هناك ($p$ قابل للقلب بترديد $q^k$: بيزو)، ومنه فحاصل القسمة $0$.
وتمرّ المجاميع المباشرة: $\dim_{A/(p)}
p^{j-1}T/p^jT = \#\{i : k_{p,i} \geq j\}$. وتحدّد هذه الأبعاد الذاتية
المجموعةَ المتعدّدة للأُسس.
\end{proof}

\begin{corollary}[الزمر الأبيلية المولَّدة بعدد منتهٍ]
\label{cor:b3:modules:abelian}
كل زمرة أبيلية مولَّدة بعدد منتهٍ من الشكل $\Z^r \times
\Z/d_1\Z\times\dots\times\Z/d_s\Z$ مع $d_1 \mid \dots \mid
d_s$،
بطريقة وحيدة. وكل زمرة أبيلية \emph{منتهية} جداءُ زمر دائرية
رتبتها قوة لعدد أولي، وحيدةٌ بوصفها مجموعة متعدّدة.\index{زمرة
أبيلية، بنيتها}
\end{corollary}

\begin{example}\label{ex:b3:modules:abelian}
توافق الزمر الأبيلية ذات الرتبة $p^n$ \emph{تجزئات} العدد $n$:
فمن أجل $p^4$ لدينا $\Z/p^4$ و $\Z/p^3\times\Z/p$ و $(\Z/p^2)^2$ و $\Z/p^2 \times (\Z/p)^2$
و $(\Z/p)^4$ --- أي خمس زمر، لأن للعدد $4$ خمس تجزئات. وتضاعف
الرتب المختلطة الأعدادَ عددًا أوليًا عددًا أوليًا (بمبرهنة الباقي
الصيني): فتوجد $5 \times 2$ زمرة أبيلية رتبتها $2^4 \cdot 3^2 = 144$.
\end{example}

\section{تطبيق: الصيغ الناظمية للتشاكلات الذاتية}

ليكن $K$ حقلًا، وليكن $V$ فضاءً متجهيًا على $K$ بُعده المنتهي $n$،
وليكن $u \in \mathcal L(V)$؛ نجعل $V$ مقاسًا على $K[X]$ بواسطة $P
\cdot x = P(u)(x)$
(\cref{ex:b3:modules:examples}). وهذا \omterm{def:b3:modules:module}{المقاس} مولَّد بعدد منتهٍ من العناصر (لأن أساسًا
على $K$ يولّده) والتوائي: إذ إن المتجهات $n+1$ التالية $x, u(x), \dots, u^n(x)$
مرتبطة على $K$ من أجل كل $x$، فتعطي كثير حدود مُبيدًا غير معدوم.

\begin{definition}\label{def:b3:modules:companion}
من أجل $P = X^m + a_{m-1}X^{m-1} + \dots + a_0$ واحديًا، تكون \emph{المصفوفة
المرافقة}\index{مصفوفة مرافقة} هي
\[
C_P = \begin{pmatrix}
0 & & & -a_0\\
1 & \ddots & & -a_1\\
& \ddots & 0 & \vdots\\
& & 1 & -a_{m-1}
\end{pmatrix} :
\]
وهي مصفوفة «الضرب في $X$» على $K[X]/(P)$ في الأساس $1, \bar X, \dots, \bar X^{m-1}$.
\end{definition}

\begin{theorem}[فروبينيوس: الصيغة الناظمية الناطقة]
\label{thm:b3:modules:frobenius}
توجد متتالية وحيدة من كثيرات الحدود الواحدية غير الثابتة $P_1
\mid P_2 \mid \dots \mid P_s$
(وتُسمّى \emph{صوامد التشابه}\index{صوامد التشابه} للتشاكل $u$)
بحيث يكون، بوصفهما مقاسين على $K[X]$،
\[
V \cong K[X]/(P_1) \oplus \cdots \oplus K[X]/(P_s):
\]
وفي أساس ملائم، يكون للتشاكل $u$ مصفوفةٌ قطرية بالكتل $\operatorname{diag}(C_{P_1}, \dots, C_{P_s})$. وعلاوة
على ذلك:
\begin{enumerate}
  \item $P_s = \mu_u$ (كثير الحدود الأدنى) و $P_1\cdots P_s =
        \chi_u$ (كثير
        الحدود المميّز)؛ وعلى وجه الخصوص $\mu_u
        \mid \chi_u$ \emph{(وهي
        مبرهنة كايلي--هاملتون مُبرهَنًا عليها من جديد)} و $\chi_u \mid \mu_u^{\,s}$،
        ومنه فإن للمقدارين $\chi_u$ و $\mu_u$ العواملَ نفسها غير القابلة
        للاختزال.
  \item يكون تشاكلان ذاتيان (أو مصفوفتان مربعتان) متشابهين إذا
        وفقط إذا كان لهما \omterm{thm:b3:modules:frobenius}{صوامد التشابه} نفسها.
\end{enumerate}
\end{theorem}

\begin{proof}
نطبّق مبرهنة البنية (\cref{thm:b3:modules:structure}) على \omterm{def:b3:rings:pidufd}{حلقة المثاليات الرئيسية}
$K[X]$: فيتفكّك \omterm{def:b3:modules:module}{المقاس} \omterm{def:b3:modules:torsion}{الالتوائي} $V$ \omterm{thm:b3:modules:smith}{بعوامل صامدة} $P_i$، منظَّمة
واحدية (لأن العناصر القابلة للقلب في $K[X]$ هي $K^\times$)؛ ولا
يظهر جزء حر (لأن $V$ التوائي). وعلى كل عامل دائري $K[X]/(P_i)$،
يكون للضرب في $X$ المصفوفةُ $C_{P_i}$ في أساس قوى $\bar X$: وبضم
الأسس نحصل على الصيغة بالكتل.

(1) مُبيد $V = \bigoplus K[X]/(P_i)$ هو $(P_1)\cap
\dots\cap(P_s) = (P_s)$ (بسلسلة القابلية للقسمة: فإن $P_s$
مضاعف مشترك، ويُبيد صنفَ $1$ في العامل الأخير بالضبط المثاليُّ
$(P_s)$): إذن $\mu_u = P_s$. أما من أجل $\chi_u$: فعلى عامل دائري
لدينا $\chi_{C_P} = P$، بالتراجع على $m = \deg P$. وبنشر $\det(XI_m - C_P)$
على السطر الأول (وعناصره $X$، ثم أصفار، ثم $a_0$ في العمود
الأخير):
\[
\det(XI_m - C_P) = X\,\det\bigl(XI_{m-1} - C_{\tilde P}\bigr)
+ (-1)^{1+m}\,a_0\,\det L ,
\]
حيث $\tilde P = X^{m-1} + a_{m-1}X^{m-2} + \dots + a_1$ (بالشكل نفسه، وبقياس أصغر بواحد) و $L$ مثلثية قطرها
$(-1, \dots, -1)$، ومنه $\det L = (-1)^{m-1}$. وبالتراجع يكون الحد الأول $X\tilde P$،
والثاني $a_0$: فالمجموع $X\tilde P + a_0 = P$ (وحالة الأساس $m=1$: $\det(X + a_0) = P$).
وتتضارب المحدّدات على الكتل: $\chi_u = \prod P_i$. وأما كايلي--هاملتون:
لدينا $\chi_u \in (\mu_u)$ لأن $P_s \mid$ كلَّ واحد منها\dots\ وبالعكس، لدينا
$P_i \mid P_s$ من أجل كل $i$، ومنه فإن $\chi_u = \prod P_i$ يقسم $P_s^{\,s} = \mu_u^s$؛ ويقسم $\mu_u =
P_s$
العنصرَ $\chi_u$ لأنه أحد عوامله.

(2) التشاكلان الذاتيان المتشابهان هما بنيتا \omterm{def:b3:modules:module}{مقاس} مقترنتان، ومنه
فلهما الصوامد نفسها (بالوحدانية في \cref{thm:b3:modules:structure})؛ وبالعكس، تعطي
الصوامد المتساوية مقاسين متماثلين على $K[X]$، وتماثل المقاسات هو
بالضبط تقابل خطي يتبادل مع التشاكلين الذاتيين: أي تشابه.
\end{proof}

\begin{corollary}[التشابه لا يتأثر بامتداد الحقل]
\label{cor:b3:modules:descent}
ليكن $K \subseteq L$ حقلين وليكن $M, N \in M_n(K)$. إذا كانت $M$ و $N$
متشابهتين على $L$، فهما متشابهتان على $K$.
\end{corollary}

\begin{proof}
تُحسب \omterm{thm:b3:modules:frobenius}{صوامد التشابه} للمصفوفة $M$ بصيغة سميث للمحدّدات الجزئية
(\cref{thm:b3:modules:smith}) مطبَّقةً على مصفوفة التقديم $XI_n - M$ على $K[X]$
--- وبالفعل، فللمقاس $V_M = K^n$ على $K[X]$ التقديمُ $XI_n - M$:
إذ إن التطبيق $K[X]^n \to V_M$ المعطى بالعلاقة $(Q_i) \mapsto \sum Q_i(M)e_i$ شاملٌ ونواته
مولَّدة بأعمدة $XI_n - M$ (بتحقّق مباشر: فبترديد تلك الأعمدة،
يختزل كل عنصر من $K[X]^n$ إلى متجهة ثابتة، والمتجهات الثابتة
تُرسَل تقابليًا؛ وتفصّل مسألة نهاية الأسبوع ذلك). والقواسم
المشتركة الكبرى لكثيرات الحدود لا تتغير بامتداد الحقل: فإذا كان
$d$ \omterm{lem:b3:rings:bezout}{القاسم المشترك الأكبر} الواحدي في $K[X]$ لعائلة $(f_j)$، أعطت
متطابقة بيزو $d = \sum u_jf_j$ حيث $u_j \in K[X]$، ومنه فإن كل قاسم مشترك
للعناصر $f_j$ \emph{في $L[X]$} يقسم $d$؛ وبما أن $d$ نفسه قاسم
مشترك، فهو \omterm{lem:b3:rings:bezout}{القاسم المشترك الأكبر} في $L[X]$ أيضًا. ومنه فإن
\omterm{thm:b3:modules:smith}{العوامل الصامدة} للمصفوفة $XI_n - M$، وهي خوارج القواسم المشتركة
الكبرى المتتالية للمحدّدات الجزئية، هي نفسها على $K$ وعلى $L$:
أي إن للمصفوفتين $M$ و $N$ \omterm{thm:b3:modules:frobenius}{صوامد التشابه} نفسها على $L$ إذا وفقط إذا كان
ذلك على $K$؛ ونخلص من \cref{thm:b3:modules:frobenius}(2).
\end{proof}

\begin{theorem}[صيغة جوردان، مستنبطةً من جديد]\label{thm:b3:modules:jordan}
لنفترض أن $\chi_u$ ينشطر على $K$ (مثلًا $K = \C$). بتطبيق تفكيك
\omterm{thm:b3:modules:structure}{القواسم الأوّلية} على $V$ (كما في برهان \cref{thm:b3:modules:structure}) بدل \omterm{thm:b3:modules:smith}{العوامل
الصامدة}:
\[
V \cong \bigoplus_{\lambda, j} K[X]\big/\bigl((X -
\lambda)^{k_{\lambda,j}}\bigr),
\]
وفي الأساس $\bigl(\overline{(X-\lambda)^{k-1}}, \dots,
\overline{(X - \lambda)}, \bar 1\bigr)$ لكل عامل، يفعل $u$ بوصفه كتلة جوردان
$J_k(\lambda)$: أي إن لكل تشاكل ذاتي ينشطر كثيرُ حدوده المميّز
أساسَ جوردان، وتكون المجموعة المتعدّدة للكتل $(\lambda, k)$
وحيدة.\index{صيغة جوردان}
\end{theorem}

\begin{proof}
\omterm{thm:b3:modules:structure}{القواسم الأوّلية} \omterm{def:b3:modules:module}{للمقاس} \omterm{def:b3:modules:torsion}{الالتوائي} $V$ هي العناصر $(X - \lambda)^k$ حيث
يمسح $X - \lambda$ العواملَ غير القابلة للاختزال للمقدار $\mu_u$ (وهو
ينشطر، لأن $\chi_u$ ينشطر ولأن لهما العوامل نفسها غير القابلة
للاختزال، \cref{thm:b3:modules:frobenius}). وفي $W = K[X]/((X-\lambda)^k)$، نضع $f_j = \overline{(X - \lambda)^{k-j}}$ من أجل $j = 1, \dots, k$:
عندئذٍ $(X - \lambda)f_j = f_{j-1}$ (مع $f_0 = 0$)، أي $u(f_j)
= \lambda f_j + f_{j-1}$: ومنه فإن مصفوفة $u$
على $(f_1, \dots,
f_k)$ هي بالضبط $J_k(\lambda)$ (بآحاد فوق القطر). ووحدانية
المجموعة المتعدّدة للقواسم الأوّلية هي \cref{thm:b3:modules:structure}.
\end{proof}

\begin{remark}\label{rem:b3:modules:overview}
صار تدرّج الصيغ الناظمية شفافًا الآن: فالصيغة \emph{الناطقة}
موجودة على كل حقل وتكشف التشابه كشفًا مطلقًا (\cref{cor:b3:modules:descent})؛
والصيغة \emph{الجوردانية} تنقيحٌ لها حين ينشطر $\chi_u$. وقد
انطوت براهين مبرهنة جوردان القائمة على إحصاء الأبعاد في السنة
الجامعية 2 تحت هذا الإطار: فكل تلك التوافيق لم تكن سوى حسابٍ في
\omterm{def:b3:rings:pidufd}{حلقة المثاليات الرئيسية} $K[X]$.
\end{remark}

\section{تمارين}

\begin{exercise}[$\star$]\label{exo:b3:modules:1}
(a) برهن على أن $\Q$ ليس \omterm{def:b3:modules:free}{مولَّدًا بعدد منتهٍ من العناصر} بوصفه
مقاسًا على $\Z$.
(b) برهن على أن $\Q$ عديم \omterm{def:b3:modules:torsion}{الالتواء} لكنه ليس \omterm{def:b3:modules:free}{حرًّا}.
(c) ولماذا لا تناقض أيٌّ من العبارتين \cref{thm:b3:modules:structure}؟
\end{exercise}

\begin{exercise}[$\star$]\label{exo:b3:modules:2}
اسرد الزمر الأبيلية ذات الرتبة $360$ إلى حدود التماثل، بصيغتَي
\omterm{thm:b3:modules:structure}{القواسم الأوّلية} \omterm{thm:b3:modules:smith}{والعوامل الصامدة} معًا. وكم زمرة أبيلية رتبتها
$p^5$؟
\end{exercise}

\begin{exercise}[$\star$]\label{exo:b3:modules:3}
احسب صيغة سميث الناظمية على $\Z$ للمصفوفتين
\[
B = \begin{pmatrix} 2 & 4\\ 6 & 8 \end{pmatrix},
\qquad
C = \begin{pmatrix} 2 & 0 & 0\\ 0 & 3 & 0\\ 0 & 0 & 12
\end{pmatrix},
\]
وعيّن الزمرتين الأبيليتين $\Z^2/B\Z^2$ و $\Z^3/C\Z^3$.
\end{exercise}

\begin{exercise}[$\star\star$]\label{exo:b3:modules:4}
لتكن $L \subseteq \Z^n$ زمرة جزئية رتبتها $n$ أساسُها أعمدة المصفوفة
$B \in M_n(\Z)$ مع $\det B \neq 0$. برهن على أن $\Z^n/L$ منتهية
وعدد عناصرها $\abs{\det B}$، وأن $\Z^n/L
\cong \prod_i \Z/d_i\Z$ حيث $d_i$ هي \omterm{thm:b3:modules:smith}{العوامل
الصامدة} للمصفوفة $B$. وأعطِ مثالًا توضيحيًا بالمصفوفة $L = \Z(2,0) + \Z(1,3)$.
\end{exercise}

\begin{exercise}[$\star\star$]\label{exo:b3:modules:5}
لتكن $A$ حلقة تامة. (a) تحقق من أن $T(M)$ \omterm{def:b3:modules:module}{مقاس} جزئي وأن $M/T(M)$
عديم \omterm{def:b3:modules:torsion}{الالتواء}. (b) برهن على أن \omterm{def:b3:rings:ideal}{المثالي} $(X, Y)$ في $K[X,Y]$،
بوصفه مقاسًا على $K[X,Y]$، عديم \omterm{def:b3:modules:torsion}{الالتواء} لكنه ليس \omterm{def:b3:modules:free}{حرًّا}: فمبرهنة
البنية تحتاج فعلًا إلى فرضية المثاليات الرئيسية.
\end{exercise}

\begin{exercise}[$\star\star$]\label{exo:b3:modules:6}
(a) برهن على أنه ليس \omterm{def:b3:modules:module}{للمقاس} $2\Z$ متمّم مباشر في \omterm{def:b3:modules:module}{المقاس} $\Z$ على $\Z$:
فالمقاسات الجزئية من \omterm{def:b3:modules:free}{مقاس حر} حرةٌ (\cref{thm:b3:modules:submodule})، لكنها ليست
بالضرورة حدودًا مباشرة. (b) برهن على أنه إذا حقّق $M \subseteq A^n$ (حيث
$A$ \omterm{def:b3:rings:pidufd}{حلقة مثاليات رئيسية}) الشرطَ: $A^n/M$ عديم \omterm{def:b3:modules:torsion}{الالتواء}، فإن $M$
\emph{يكون} حدًّا مباشرًا.
\end{exercise}

\begin{exercise}[$\star\star$]\label{exo:b3:modules:7}
حل في $\Z^2$ الجملة
\[
\begin{cases}
2x + 4y \equiv b_1 \pmod{20},\\
6x + 8y \equiv b_2 \pmod{20},
\end{cases}
\]
وعيّن الأزواج $(b_1, b_2)$ التي توجد لها حلول، باستعمال صيغة
سميث في \cref{exo:b3:modules:3} (بتغييرات متغيرات قابلة للقلب في الطرفين).
\end{exercise}

\begin{exercise}[$\star\star$]\label{exo:b3:modules:8}
(a) عيّن جميع \omterm{thm:b3:modules:frobenius}{صوامد التشابه} وصيغ جوردان الممكنة لمصفوفة معدومة
القوة من القياس $4 \times 4$، مرتّبةً حسب تجزئة العدد $4$ التي
تحقّقها.
(b) أعطِ مصفوفتين عقديتين من القياس $4\times4$ لهما كثيرا الحدود
المميّز \emph{و} الأدنى نفساهما وليستا متشابهتين، وبرهن على أن
ذلك لا يمكن أن يحدث من أجل $n \leq 3$.
\end{exercise}

\begin{exercise}[$\star\star\star$]\label{exo:b3:modules:9}
ليكن $u \in \mathcal L(V)$ مع $\dim V = n$. برهن على تكافؤ ما يلي: (أ) يكون
$V$ مقاسًا دائريًا على $K[X]$ (أي يوجد $x$ يحقق $V = K[u]x$،
ويُسمّى \emph{متجهة دائرية})؛ (ب) $\mu_u = \chi_u$؛ (ج)
$s = 1$ في \cref{thm:b3:modules:frobenius}. واستنتج أن لكل \omterm{def:b3:modules:companion}{مصفوفة مرافقة} متجهةً
دائرية، وعيّن متى يكون لمصفوفة قطرية متجهةٌ دائرية.
\end{exercise}

\begin{exercise}[$\star\star\star$]\label{exo:b3:modules:10}
من أجل $M \in M_n(\Z)$ منظورًا إليها بوصفها تشاكلًا ذاتيًا
\omterm{def:b3:modules:module}{للمقاس} $\Z^n$، برهن على صيغة الدليل: إذا كان $\det M \ne 0$ فإن
$[\Z^n : M\Z^n] = \abs{\det M}$، واستنتج أن $M \in GL_n(\Z)$ إذا وفقط إذا كان
$\det M = \pm 1$. تطبيق: للزمرة $\Z^2$ عدد $\sigma_1(m) = \sum_{d \mid m} d$ من الزمر
الجزئية ذات الدليل $m$ بالضبط. \emph{(أحصِ المصفوفات في صيغة
هيرميت $\bigl(\begin{smallmatrix}
a & b\\ 0 & d\end{smallmatrix}\bigr)$ مع $ad = m$ و $0 \leq b <
d$.)}
\end{exercise}

\begin{exercise}[$\star\star$]\label{exo:b3:modules:11}
(المعادلات $x^k = e$ في الزمر الأبيلية) لتكن $G$ زمرة أبيلية
منتهية عواملها الصامدة $d_1 \mid d_2 \mid \dots
\mid d_s$.
(a) برهن على أنه من أجل كل $k \geq 1$،
\[
\#\{x \in G : x^k = e\} \;=\; \prod_{i=1}^{s}\gcd(k, d_i) .
\]
(b) استنتج: تكون زمرة أبيلية منتهية دائرية إذا وفقط إذا كان
للمعادلة $x^k = e$ عدد حلول لا يتجاوز $k$ من أجل كل $k$.
(c) استعد دائرية الزمر الجزئية المنتهية من $K^\times$ (حيث $K$
حقل، \cref{ch:b3:galois}): ولماذا يضمن كثير الحدود $X^k - 1$ محكَّ (b)؟
\end{exercise}

\begin{exercise}[$\star\star\star$]\label{exo:b3:modules:12}
(المصفوفات الأوّلية تولّد) (a) برهن على أن $M \in M_n(\Z)$ قابلة
للقلب في $M_n(\Z)$ إذا وفقط إذا كان $\det M = \pm1$.
(b) برهن على أن $SL_2(\Z)$ مولَّدة بالمصفوفتين الأوّليتين
$E = \bigl(\begin{smallmatrix}1 & 1\\ 0 &
1\end{smallmatrix}\bigr)$ و $F = \bigl(\begin{smallmatrix}1 & 0\\ 1 &
1\end{smallmatrix}\bigr)$.
\emph{(شغّل خوارزمية إقليدس على العمود الأول للمصفوفة $M
\in SL_2(\Z)$
بالضرب من اليسار في قوى $E, F$، وصولًا إلى $\pm\bigl(\begin{smallmatrix}1 & *\\ 0 &
1\end{smallmatrix}\bigr)$؛ ثم أنهِ
يدويًا --- ولاحظ $-I = (EF^{-1}E)^2$.)}
(c) اشرح الصلة باختزال سميث: فعلى $\Z$، تكفي العمليات على الأسطر
والأعمدة ذات المحدد $1$ للتقطير، إلى حدود الإشارات.
\end{exercise}

\section{مسألة: المُبادِل والمُبادِل المضاعف}

\begin{problem}[{مسألة نهاية الأسبوع --- الصيغة الناطقة،
والمُبادِل، والمُبادِل المضاعف}]\label{pb:b3:modules:1}
ليكن $K$ حقلًا، وليكن $V$ فضاءً متجهيًا على $K$ بُعده $n \geq
1$،
وليكن $u \in \mathcal L(V)$. ندرس \emph{المُبادِل}
\[
\mathcal C(u) = \{v \in \mathcal L(V) : uv = vu\},
\]
وهو جبر جزئي من $\mathcal L(V)$ يحتوي على $K[u] = \{P(u) : P \in
K[X]\}$، ونبرهن على
صيغة فروبينيوس للبُعد وعلى مبرهنة المُبادِل المضاعف: $\mathcal C(\mathcal C(u)) = K[u]$. وفي
كل ما يلي، $V$ هو \omterm{def:b3:modules:module}{المقاس} على $K[X]$ المعرَّف بالتشاكل $u$، وعوامله
الصامدة $P_1 \mid \cdots \mid P_s$ وتفكيكه الدائري $V =
\bigoplus_{i=1}^s V_i$، $V_i = K[u]\,x_i \cong K[X]/(P_i)$، $n_i
= \deg P_i$
(\cref{thm:b3:modules:frobenius}).

\medskip
\noindent\textbf{الجزء الأول --- مصفوفة التقديم $XI - M$،
وتمارين تمهيدية.}
\begin{enumerate}
  \item لتكن $M \in M_n(K)$ وليرسل $\varphi \colon K[X]^n \to V_M
        = K^n$ العنصرَ $(Q_1, \dots, Q_n)$ إلى
        $\sum_i Q_i(M)e_i$. برهن على أن $\varphi$ تشاكل شامل لمقاسات على
        $K[X]$ وأن كل عمود من $XI_n - M$ يقع في $\ker\varphi$.
  \item برهن على أن كل عنصر من $K[X]^n$ موافقٌ، بترديد أعمدة
        $XI_n - M$، لمتجهة ثابتة \emph{(اختزل الدرجات باستعمال
        $Xe_i \equiv \sum_j m_{ji}e_j$)}، واستنتج $\ker\varphi = (XI_n - M)\,K[X]^n$: أي إن \omterm{def:b3:modules:module}{للمقاس} $V_M$ مصفوفةَ
        التقديم $XI_n - M$. واستعد نقطة انطلاق \cref{cor:b3:modules:descent}: أي
        إن \omterm{thm:b3:modules:frobenius}{صوامد التشابه} للمصفوفة $M$ هي \omterm{thm:b3:modules:smith}{العوامل الصامدة} غير
        القابلة للقلب للمصفوفة $XI_n - M$.
  \item احسب \omterm{thm:b3:modules:frobenius}{صوامد التشابه} لكلٍّ من: مصفوفة سلَّمية $\lambda I_n$؛
        ومصفوفة قطرية عناصر قطرها متمايزة؛ وكتلة جوردان
        $J_n(0)$ من القياس $n \times n$؛ والمصفوفة $\operatorname{diag}(J_2(0), J_1(0))$ من
        أجل $n = 3$.
  \item برهن على أن $\dim K[u] = \deg \mu_u = n_s$.
\end{enumerate}

\medskip
\noindent\textbf{الجزء الثاني --- التشاكلات بين المقاسات
الدائرية.}
\begin{enumerate}[resume]
  \item ليكن $P, Q$ واحديين غير ثابتين. برهن على أن كل تشاكل
        $f \colon K[X]/(P) \to K[X]/(Q)$ على $K[X]$ يتحدّد بالقيمة $f(\bar 1)$، وأن
        $\bar c \in K[X]/(Q)$ يصلح أن يكون $f(\bar 1)$ إذا وفقط إذا كان
        $P\bar c = 0$ في $K[X]/(Q)$.
  \item استنتج
        \[
        \operatorname{Hom}_{K[X]}\bigl(K[X]/(P),\, K[X]/(Q)\bigr)
        \;\cong\; K[X]\big/\bigl(\gcd(P, Q)\bigr),
        \]
        وبُعده $\deg \gcd(P, Q)$ على $K$. \emph{(برهن على أن الحلول $\bar c$ للمعادلة
        $P\bar c = 0$ في $K[X]/(Q)$ تشكّل \omterm{def:b3:modules:module}{المقاس} الجزئي الدائري
        المولَّد بالعنصر $\overline{Q/\gcd(P,Q)}$.)}
  \item برهن على \emph{صيغة فروبينيوس}:
        \[
        \dim_K \mathcal C(u) = \sum_{i,j=1}^{s} \deg\gcd(P_i,
        P_j) = \sum_{i=1}^{s} (2s - 2i + 1)\, n_i .
        \]
        \emph{(التشاكل $v$ المتبادل مع $u$ هو بالضبط تشاكل ذاتي
        \omterm{def:b3:modules:module}{للمقاس} $V$ على $K[X]$؛ فكّك $\operatorname{End}(\bigoplus_i V_i)$ إلى مصفوفات تشاكلات
        $V_j \to V_i$ واستعمل سلسلة القابلية للقسمة.)}
  \item استنتج $\dim \mathcal C(u) \geq n$، مع التساوي إذا وفقط إذا كان $u$ دائريًا
        ($s = 1$)، واحسب $\dim\mathcal C(u)$ من أجل $u = \lambda\,\mathrm{id}$: وهما طرفا
        الصيغة.
  \item تحقق من صيغة فروبينيوس مباشرةً من أجل $\operatorname{diag}(J_2(0), J_1(0))$ بحساب
        المُبادِل صراحةً على صورة مصفوفات من القياس $3\times3$.
\end{enumerate}

\medskip
\noindent\textbf{الجزء الثالث --- مبرهنة المُبادِل المضاعف.}
ليكن $w \in \mathcal C(\mathcal C(u))$؛ نبرهن على أن $w \in K[u]$.
\begin{enumerate}[resume]
  \item برهن على $K[u] \subseteq \mathcal C(\mathcal C(u))$، وعلى أن كل $w \in \mathcal C(\mathcal C(u))$ يتبادل مع $u$ ---
        ومنه فإن الاحتواء المطلوب برهانه، $\mathcal C(\mathcal C(u)) \subseteq K[u]$، تدقيقٌ حقيقي
        للاحتواء $w \in \mathcal C(u)$.
  \item لنفترض أولًا أن $u$ \emph{دائري}، أي $V = K[u]x$. برهن
        مباشرةً على أن $\mathcal C(u) = K[u]$ \emph{(قيّم تشاكلًا متبادلًا $v$
        عند $x$: فيكون $v(x) = P(u)x$ من أجل $P$ ما، ثم قارن $v$
        بالتشاكل $P(u)$ على الأساس $u^k x$)}، واخلص إلى المبرهنة في هذه
        الحالة.
  \item لنعد إلى الحالة العامة. من أجل كل $i$، ليكن $\pi_i\colon
        V \to V_i$
        الإسقاطَ على امتداد الحدود الأخرى. برهن على $\pi_i \in \mathcal C(u)$،
        واستنتج أن $w$ يحفظ كل $V_i$ ويتبادل مع $u_i =
        u\restriction_{V_i}$؛ ثم اخلص
        بواسطة السؤال 11 مطبَّقًا على $u_i$ الدائري: أي توجد
        كثيرات حدود $Q_i$ تحقق $w\restriction_{V_i} = Q_i(u)\restriction_{V_i}$.
  \item يبقى أن نلصق كثيرات الحدود $Q_i$ في كثير حدود واحد. من
        أجل $i
        \leq j$ (ومنه $P_i \mid P_j$)، برهن على أن $\eta_{ij} \colon
        V_j \to V_i$،
        $R(u)x_j \mapsto R(u)x_i$، تشاكلٌ معرَّف تعريفًا سليمًا على $K[X]$
        \emph{(وما يجب التحقق منه هو أن $R(u)x_j = 0$ يستلزم
        $R(u)x_i = 0$)}، وأن $\tilde\eta_{ij} = \eta_{ij}\circ\pi_j$، ممدَّدًا بالقيمة $0$ على الحدود
        الأخرى، يقع في $\mathcal C(u)$.
  \item باستعمال $w\tilde\eta_{ij} = \tilde\eta_{ij}w$، برهن على $Q_i
        \equiv Q_j \pmod{P_i}$ من أجل $i \leq j$.
        واستنتج أن $Q =
        Q_s$ يحقق $Q \equiv Q_i \pmod {P_i}$ من أجل كل $i$، ومنه
        $w = Q(u)$ على كل $V_i$، ومنه على $V$:
        \[
        \boxed{\ \mathcal C(\mathcal C(u)) = K[u].\ }
        \]
  \item (خاتمة) استنتج من المبرهنة: إذا كان $v$ يتبادل مع
        \emph{كل} مصفوفة تتبادل مع $u$، وكان $u$ دائريًا، فإن $v$
        كثير حدود في $u$؛ وأعطِ مثالًا يبيّن أن $\mathcal C(u) = K[u]$ يخفق من
        أجل $u =
        \mathrm{id}$ حيث $n \geq 2$ --- وأين تدخل الدائرية
        بالضبط؟
\end{enumerate}

\medskip
\noindent\textbf{الجزء الرابع --- عوائد \omterm{thm:b3:modules:frobenius}{صوامد التشابه}.} الصيغة
الناظمية الناطقة آلة؛ وهذه خمسة من مخرجاتها الكلاسيكية.
\begin{enumerate}[resume]
  \item (المنقولة) برهن على أن كل $M \in M_n(K)$ تشابه منقولتها
        ${}^tM$. \emph{(فالعمليات التي تحوّل $XI - M$ إلى صيغة
        سميث، منقولةً، تحوّل $XI - {}^tM$ إلى صيغة سميث نفسها:
        أي صوامد تشابه متساوية.)}
  \item (نزول التشابه) ليكن $K \subseteq L$ امتداد حقول وليكن
        $M, N \in M_n(K)$. برهن على أنه إذا كانت $M$ و $N$ متشابهتين على
        $L$، فهما متشابهتان على $K$. \emph{(فصيغة سميث للمصفوفة
        $XI - M$ المحسوبة في $K[X]$ تبقى صيغة سميث في $L[X]$ ---
        ولماذا لا تتغير \omterm{thm:b3:modules:smith}{العوامل الصامدة}؟)} ونتيجة جديرة بالحفظ:
        كل مصفوفتين حقيقيتين مقترنتين في $GL_n(\C)$ مقترنتان في
        $GL_n(\R)$.
  \item (تصنيف معدومات القوة) ليكن $u$ معدوم القوة. برهن على أن
        عدد الكتل ذات القياس $\geq k$ في تفكيكه إلى كتل جوردان
        معدومة القوة يساوي $\operatorname{rk}u^{k-1} - \operatorname{rk}u^k$، واستنتج: أن أصناف معدومات
        القوة في $M_n(K)$، من أجل أي حقل $K$، في تقابل مع تجزئات
        العدد $n$. وكم صنفًا من معدومات القوة في $M_5(K)$؟
  \item (زوج ملموس) عيّن \omterm{thm:b3:modules:frobenius}{صوامد التشابه} للاشتقاق $D\colon P \mapsto P'$ الفاعل
        على فضاء كثيرات الحدود ذات الدرجة $< n$ المرموز إليه
        بالرمز $K_{n-1}[X]$: (a) من أجل $K =
        \Q$؛ (b) من أجل $K = \mathbb F_p$
        حيث $p < n$ \emph{(ففي المميّز $p$ لدينا $(X^p)' = 0$:
        احسب $\ker D^k$ واستعمل السؤال 18)}.
  \item (\omterm{ex:b3:groups:actions}{أصناف الاقتران} في $GL_2(\mathbb F_q)$) باستعمال \omterm{thm:b3:modules:smith}{العوامل الصامدة}،
        برهن على أن كل صنف من $GL_2(\mathbb F_q)$ من نمط واحد بالضبط من
        الأنماط الأربعة: مركزي $aI$؛ أو قابل للتقطير بقيمتين
        ذاتيتين متمايزتين $a \neq b$ في $\mathbb F_q^\times$؛ أو غير نصف بسيط
        وكثيرُ حدوده الأدنى $(X - a)^2$؛ أو دائري وكثيرُ حدوده
        المميّز غير قابل للاختزال.
  \item أحصِ أصناف كل نمط واخلص إلى: أن للزمرة $GL_2(\mathbb F_q)$ عدد $q^2 - 1$
        من \omterm{ex:b3:groups:actions}{أصناف الاقتران} بالضبط. \emph{(أحصِ كثيرات الحدود
        الواحدية غير القابلة للاختزال من الدرجة الثانية على
        $\mathbb F_q$؛ والأزواج غير المرتّبة $\{a, b\}$؛ وتذكّر
        أن القابلية للقلب تقيّد الحدود الثابتة.)}
  \item (الدائرية هي الحالة العامة) برهن على أن $M \in M_2(\mathbb F_q)$ تخفق في
        أن تكون دائرية إذا وفقط إذا كانت $M$ سلَّمية، واستنتج أن
        مصفوفة عشوائية منتظمة من القياس $2\times2$ على
        $\mathbb F_q$ تكون دائرية باحتمال $1 - q^{-3}$. واذكر
        الحدس المماثل من أجل $M_n$ ومن أجل $q$ كبير (دون برهان):
        فالمصفوفات غير الدائرية نادرة --- ولهذا لم يحتج الجزء
        الثالث من \cref{pb:b3:modules:1} إلى عمل حقيقي إلا بعد تجاوز الحالة
        العامة.
\end{enumerate}

\medskip
\noindent\textbf{الجزء الخامس --- تكملات.}
\begin{enumerate}[resume]
  \item (مركز المُبادِل) برهن على أن مركز الجبر $\mathcal C(u)$
        هو بالضبط $K[u]$ \emph{(اجمع بين احتوائَي الجزء الثالث)}.
        واستنتج أن $\mathcal C(u)$ تبديلي إذا وفقط إذا كان $u$
        دائريًا --- فتستعيد حالة التساوي في السؤال 8 بطريق بنيوي
        محض.
  \item (أي الأبعاد تتحقق؟) استنتج من صيغة فروبينيوس أن $\dim\mathcal C(u) \equiv n \pmod 2$
        من أجل كل $u$. ثم عيّن بدقة مجموعة القيم التي يأخذها
        $\dim \mathcal C(u)$ حين يمسح $u$ الفضاءَ $\mathcal L(V)$ مع
        $\dim V = 4$: برهن على أنها $\{4, 6,
        8, 10, 16\}$ \emph{(اسرد متتاليات
        الدرجات $n_1
        \leq \dots \leq n_s$ التي مجموعها $4$ وحقّق كلًّا منها بمصفوفة
        معدومة القوة)}. وعلى وجه الخصوص لا يتحقق العددان $12$
        و $14$ رغم أن زوجيتهما صحيحة: فقيد الزوجية لازم وغير
        كافٍ.
  \item (\omterm{cor:b3:groups:classeq}{معادلة الأصناف} في $GL_2(\mathbb F_3)$) من أجل $q = 3$، احسب عدد
        عناصر كل \omterm{ex:b3:groups:actions}{صنف اقتران} من أصناف السؤال 20 بمبرهنة المدار
        والمثبِّت: فمُرَكِّز مصفوفة $M$ \emph{دائرية} في $GL_2(\mathbb F_q)$
        هو زمرة العناصر القابلة للقلب من $K[M]$ (السؤال 11).
        عيّن $K[M]$ في الأنماط الثلاثة غير المركزية، واسرد
        كثيرات الحدود الواحدية الثلاثة غير القابلة للاختزال من
        الدرجة الثانية على $\mathbb F_3$، وتحقق من \omterm{cor:b3:groups:classeq}{معادلة
        الأصناف}
        \[
        48 = \abs{GL_2(\mathbb F_3)}
        = 2\cdot1 + 1\cdot12 + 2\cdot8 + 3\cdot6,
        \]
        بعدد $2 + 1 + 2 + 3 = 8 = q^2 - 1$ من الأصناف، كما تنبّأ السؤال 21.
\end{enumerate}
\end{problem}
